\section{スタイルガイドとの対応}\label{sec:guidemap}
この章は，社会学評論スタイルガイド（第3版）の各節が，このスタイルのどこで扱われ，どの回帰テストで確かめられているかの対応表です。スタイルガイドの文章は引用せず，節番号と要点だけを示します。

「テスト」の欄の C01 などは，\file{tests/guide.expected.txt} の文献注のケース番号です。番号のないものは，同じファイルの注釈（\texttt{\#} で始まる行）に節番号があります。期待値はスタイルガイドの例どおりに書いたもので，実際の出力に合わせて書き換えることはしません。テストの実行方法は，リポジトリの \file{tests/README.md} にあります。

\subsection{対応表}
\begingroup\small
\setlength{\tabcolsep}{4pt}
\begin{longtable}{@{}L{.11\textwidth}L{.31\textwidth}L{.17\textwidth}L{.26\textwidth}@{}}
\toprule
\textbf{ガイド} & \textbf{内容} & \textbf{この文書} & \textbf{テスト} \\
\midrule
\endfirsthead
\toprule
\textbf{ガイド} & \textbf{内容} & \textbf{この文書} & \textbf{テスト} \\
\midrule
\endhead
\bottomrule
\endfoot
3.1 & 研究者名の表記（初出はフルネーム，2回目以降は姓のみ）$^{*}$ & \ref{sec:firstfull} 節 & C01，C05，C06，C41〜C44，C46〜C48，C53，C55，C56 \\
3.2 & 文献を示す割注（ページ，共著，編者，訳書，再録，同姓の著者，3名以上） & \ref{sec:cite} 節，\ref{sec:citetrans} 節 & C02〜C32，C35〜C40，C54，C57，C58 \\
3.3 & 参照（著者名の後の文献注） & \ref{sec:cite} 節 & C05，C06 \\
3.4〜3.6 & 短い引用・長い引用・フィールドワーク資料 & 対象外（本文の書き方） & ―― \\
3.7 & 新聞記事などからの引用 & \ref{sec:news} 節 & C34，C45，C66〜C69 \\
3.8.1 & ウェブ文書の文献注（段落指定） & \ref{sec:cite} 節 & C33 \\
3.8.2 & ウェブ文書を論文で使う場合の注意 & 対象外 & ―― \\
4.1 & 文献リストの形式（区切り，ぶら下げ，姓の順，同一著者の並び，a・b）$^{*}$ & \ref{sec:sorting} 節 & C59，C60，注釈「4.1 同一著者」，「4.4(1) 同一著者・同一年」 \\
4.2(1)〜(3) & 欧文の単著・共著・編書 & \ref{sec:book} 節 & 注釈「4.2(2) 3人以上の共著」ほか \\
4.2(4)，4.4(4) & 編書論文 & \ref{sec:incollection} 節 & 注釈「4.4(4)」，「欧文の編書論文でも，所収書の版を出す」 \\
4.2(5)，4.4(5) & 雑誌論文 & \ref{sec:article} 節 & 注釈「4.4(5) 号のみの雑誌」 \\
4.2(6)，4.4(6) & 翻訳書・英訳文献$^{*}$ & \ref{sec:trans} 節 & C19，C20，C26，C37，C38，C50，C57，C58，注釈「4.2(6)」，「4.4(6)」 \\
4.2(7) & 初版の出版年（再刊）$^{*}$ & \ref{sec:trans} 節 & C23〜C25，C27，C49，注釈「4.2(7)」 \\
4.2(8)，4.4(7) & 初出誌の掲載誌と再録編著書をともに示す$^{*}$ & \ref{sec:trans} 節 & 注釈「4.2(8)」，「4.4(7)」 \\
4.3 & その他の外国語文献 & 対象外（投稿者の責任） & ―― \\
4.4(1)〜(3) & 邦文の単著・共著・編書（版は書名に含める）$^{*}$ & \ref{sec:book} 節 & 注釈「4.4(1)」，C59，C60 \\
4.4(8) & 調査報告書$^{*}$ & \ref{sec:report} 節 & C61，C62，注釈「4.4(8)」 \\
4.4(9) & 政府刊行物（発行元の自動省略）$^{*}$ & \ref{sec:report} 節 & 注釈「4.4(9)」，「書籍：発行元が著者と同じなら省く」 \\
4.4(10) & 修士論文・学会報告原稿$^{*}$ & \ref{sec:thesis} 節 & 注釈「4.4(10)」 \\
4.5.1 & 電子書籍$^{*}$ & \ref{sec:ebook} 節 & C63〜C65，注釈「4.5.1」 \\
4.5.2 & 電子ジャーナル（取得日，URL，データベース名） & \ref{sec:article} 節 & 注釈「4.5.2.1」，「4.5.2.2」 \\
4.5.3 & ウェブページ・ブログ$^{*}$ & \ref{sec:web} 節 & 注釈「4.5.3」 \\
4.5.4.1 & オンラインデータベース & \ref{sec:web} 節 & 注釈「4.5.4.1」 \\
4.5.4.2 & コンピュータファイル（MRDF）$^{*}$ & \ref{sec:dataset} 節 & 注釈「4.5.4.2」 \\
4.5.5 & URL について & \ref{sec:cautions} 節 & 注釈「4.5.3」ほか \\
\end{longtable}
\endgroup

$^{*}$ は，スタイルガイドに明確な規定がなく，スタイルの側で解釈して決めた点を含む節です。次の節に挙げます。

\subsection{解釈で決めた点}
スタイルガイドの例や文面から一意に決まらない点は，次のように決めました。決めた経緯は，リポジトリの \file{docs/DECISIONS.md} に日付つきで記録しています。意図と違う出力になる場合は，この節の「決めたこと」を確かめてください。

\begin{description}
\item[3.1 初出の判定]
同じ人物の別の文献でも，一度出た人物は姓のみにします。判定は人物単位です（共著は名前の並び全体で1単位）。割注（\cmd{parencite}）の中の名前は，言及に数えません。\opt{firstfull=false} で，引用キーごとの判定に切り替えられます（\ref{sec:firstfull} 節）。
\item[3.2 3名以上]
割注でも本文中の引用（\cmd{textcite}）でも，「ほか」「et al.」にします。
\item[3.2 訳書のページ]
原書と訳書の両方のページを示すときは，\texttt{\textbackslash parencite[78=86]\{訳書\}} と書きます（\ref{sec:citetrans} 節）。
\item[4.1 並び順・a・b]
同一著者の並びは，単著，単独の編書，共著，共編書の順で，各カテゴリーの中は出版年順です（訳書は原典の年を優先します）。a，b は，同一著者・同一カテゴリー・同一年だけに付けます（「小田 1991」と「小田編 1991」には付けません）。
\item[4.2(6)，4.4(6) 原典の二重掲載]
訳書が引用されたときは，原典を単独の項目にせず，原典の項目に訳書を添えます。原典だけが引用されたときは，原典を単独で出します（\ref{sec:related} 節，\ref{sec:trans} 節）。
\item[4.2(7)，4.2(8)，4.4(7) 再刊と再録の使い分け]
原典が書籍で，書名・副題が同じなら，再刊として「[初版年] 手元の版の年」の形にします。原典が論文，または書名が変わった再刊は，「（再録：……）」「Reprinted in:」の形です。同じ原典の再録と翻訳をともに引用したときは，1項目にまとめ，添える項目を年順に並べます。
\item[4.4(1) 版]
和文の書籍の版は，書名の後に半角スペースで区切って出します（『社会調査入門 第2版』）。
\item[4.4(8) 報告書の先頭の名前]
\bibfield{author}，\bibfield{editor}，\bibfield{institution} の順に，あるものを先頭に出します。\bibfield{editor} を先頭にしたときは，研究機関名を種類・番号の後に出します。
\item[4.4(9) 発行元の自動省略]
発行元が著者・編者・機関名と同じとき，または「国立印刷局」（旧・大蔵省印刷局）のときは，省きます。bib には省かずに書きます。
\item[4.4(10) 学会報告]
和文の学会報告で \bibfield{type} を書かなかったときは，「報告原稿」を補いません。欧文の博士論文（\texttt{phdthesis}）は，ASA の形（PhD dissertation，括弧なし）で出します。
\item[4.5.1 電子書籍]
紙媒体の情報は \bibfield{origdate}＋\bibfield{origpublisher} に書き，フォーマットは \bibfield{note} に書きます。文献注には電子書籍の年だけを出します。
\item[4.5.3 ブログの日付]
\bibfield{date} に日まで書いてあれば出します。
\item[4.5.4.2 MRDF の提供機関]
作成者と同じでも省きません（スタイルガイドの例が省いていないため）。
\item[欧文の雑誌論文の出版社・叢書]
欧文の雑誌論文では出版社を出しません（和文では出します）。叢書は和文だけ書名の前に出し，欧文では出しません。
\item[カタカナ書きの外国人名]
「名・姓」の形で出し，共著の区切りは全角スラッシュにします（杉本良夫／ロス・マオア）。文献注にも同じ規則を適用します。
\item[欧文タイトルの大文字化]
データ側の責任です。自動で変換しません（\ref{sec:cautions} 節）。
\end{description}
